% =============================================================================
% profile_report.tex
% =============================================================================
%
% Phase 0 Hardware Profiling Report
% hallmhd-dlra-tda  research programme
%
% Records the measured FFT, SVD, and memory budget results from
% profiling/phase0_results.json (run 2026-05-17 19:26:41).
%
% Usage:
%   pdflatex profiling/profile_report.tex
% =============================================================================

\documentclass[11pt, a4paper]{article}

\usepackage{amsmath, amssymb}
\usepackage{geometry}
\usepackage{booktabs}
\usepackage{hyperref}
\usepackage{xcolor}
\usepackage{microtype}
\usepackage{tcolorbox}
\usepackage{pgfplots}
\usepackage{listings}
\usepackage{siunitx}

\pgfplotsset{compat=1.18}
\geometry{margin=2.5cm}

\hypersetup{colorlinks=true, linkcolor=blue!70!black, urlcolor=blue!70!black}

\tcbuselibrary{breakable}
\newtcolorbox{resultbox}[1]{
  colback=blue!5, colframe=blue!55!black,
  title={\textbf{#1}}, breakable, fonttitle=\bfseries\small,
}
\newtcolorbox{passbox}{
  colback=green!6, colframe=green!50!black, breakable,
}
\newtcolorbox{notebox}{
  colback=yellow!8, colframe=orange!60!black, breakable,
}

\lstset{basicstyle=\ttfamily\small, breaklines=true,
        frame=single, backgroundcolor=\color{gray!7}}

% =============================================================================
\begin{document}
% =============================================================================

\title{%
  \textbf{Phase 0 Hardware Profiling Report}\\[4pt]
  \large \texttt{hallmhd-dlra-tda} research programme\\[6pt]
  \normalsize Decision~D4: Metal/CPU throughput characterisation
}

\author{HPC / Apple Silicon Architect}

\date{Profiling run: 2026-05-17 19:26:41 (UTC+9)\\
Platform: Apple Mac Mini M1, 8\,GB unified RAM, Python~3.13.9\\
NumPy~2.3.5 $\cdot$ JAX~0.10.0 (CPU backend) $\cdot$ Apple Accelerate}

\maketitle

% =============================================================================
\begin{passbox}
\textbf{Phase 0 verdict: ALL PASS --- proceed to Phase~1.}
\smallskip

\begin{tabular}{llll}
D3 Memory budget & peak $= 0.013\,\text{GB}$ (R=32) & $\leq 13\,\text{GB}$ & \checkmark \\
D4a FFT (128³)   & Hero Run $\approx 0.41\,\text{days}$ & $\leq 5\,\text{days}$  & \checkmark \\
D4b SVD R=32     & $1.10\,\text{ms}$ per call         & $\leq 10\,\text{ms}$   & \checkmark \\
\end{tabular}
\end{passbox}

\tableofcontents
\newpage

% =============================================================================
\section{Environment}
\label{sec:env}
% =============================================================================

\begin{center}
\begin{tabular}{ll}
\toprule
Item & Value \\
\midrule
Hardware & Apple Mac Mini, Apple M1 chip \\
Unified RAM & 8\,GB \\
JAX backend & CPU (Metal backend disabled: \texttt{JAX\_PLATFORMS=cpu}) \\
JAX version & 0.10.0 \\
NumPy version & 2.3.5 \\
Python version & 3.13.9 \\
BLAS/LAPACK & Apple Accelerate (via NumPy linkage) \\
Profiling script & \texttt{scripts/run\_phase0\_profiling.sh} \\
Results file & \texttt{profiling/phase0\_results.json} \\
Elapsed time & 4.9\,s \\
\bottomrule
\end{tabular}
\end{center}

\begin{notebox}
\textbf{Note on Metal backend.}
\texttt{jax-metal} 0.10.0 raised \texttt{UNIMPLEMENTED: default\_memory\_space
is not supported} when attempting to allocate arrays on the M1 GPU.
This is a version incompatibility between JAX 0.10.0 and the installed
\texttt{jax-metal} plugin.  All profiling therefore ran on the CPU path,
which is the conservative Hero~Run planning baseline specified in the
proposal (\S8.3).  Metal GPU timings will be collected on the M4~16\,GB
machine in Phase~1 and recorded as supplementary data.
\end{notebox}

\begin{notebox}
\textbf{Note on JAX float64.}
JAX defaulted to float32 for the JAX-path FFT despite \texttt{dtype=float64}
being requested.  The NumPy path (which is the planning baseline) ran in
float64 throughout.  \textbf{Before Phase~1, add
\texttt{jax.config.update("jax\_enable\_x64", True)} to all simulation
entry points.}
\end{notebox}

% =============================================================================
\section{3D Real FFT Throughput (Step 1/4 and 2/4)}
\label{sec:fft}
% =============================================================================

\subsection{Measurement methodology}

The profiler measures the wall time of one forward + inverse 3D real FFT
roundtrip (\texttt{np.fft.irfftn(np.fft.rfftn(x))}) over
$N_\text{warmup}=5$ discarded trials and $N_\text{trials}=20$ measured
trials, reporting mean~$\pm$~std.  The ETDRK4 timestep estimate assumes
15 FFT roundtrips per step (conservative; actual count depends on
implementation but 15 is the standard estimate for RK4-class methods
with spectral dealiasing).

\subsection{Results}

\begin{center}
\begin{tabular}{lrrrrr}
\toprule
Backend & Resolution &
  Mean (ms) & Std (ms) & ETDRK4 step (ms) \\
\midrule
JAX CPU float32$^\dagger$ &
  $64\times64\times32$    & 0.29 & 0.04 &  4.4 \\
JAX CPU float32$^\dagger$ &
  $128\times128\times64$  & 2.22 & 0.27 & 33.3 \\
JAX CPU float32$^\dagger$ &
  $256\times256\times128$ & 15.33 & 0.19 & 229.9 \\
\midrule
\textbf{NumPy CPU float64} &
  $64\times64\times32$    & 1.53  & 0.22  & 22.9 \\
\textbf{NumPy CPU float64} &
  $128\times128\times64$  & 15.73 & 7.06  & 236.0 \\
\textbf{NumPy CPU float64} &
  $256\times256\times128$ & 130.03 & 2.46 & 1950.4 \\
\bottomrule
\multicolumn{5}{l}{%
  ${}^\dagger$ JAX silently downcast to float32 (see \S\ref{sec:env}).
  NumPy float64 is the planning baseline.} \\
\end{tabular}
\end{center}

\subsection{Hero Run wall-time estimates (150,000 steps)}

\begin{center}
\begin{tabular}{lrrll}
\toprule
Resolution & ETDRK4 (ms) & Hero Run (days) &
  Ceiling ($\leq 5$ days) \\
\midrule
$64\times64\times32$    &   22.9 & 0.04 & \checkmark \\
$128\times128\times64$  &  236.0 & 0.41 & \checkmark \\
$256\times256\times128$ & 1950.4 & 3.39 & \checkmark \\
\bottomrule
\end{tabular}
\end{center}

\begin{passbox}
\textbf{D4a: PASS.}
The baseline resolution $128\times128\times64$ gives a Hero Run estimate of
$\approx 0.41$ days (CPU NumPy float64), well within the 5-day ceiling.
Even $256\times256\times128$ (full-rank, CPU) is within ceiling at 3.39 days.
On the M4 Metal GPU at float32, timings will be substantially faster.
\end{passbox}

% =============================================================================
\section{Randomised SVD Throughput (Step 3/4)}
\label{sec:svd}
% =============================================================================

\subsection{Measurement methodology}

The Halko--Martinsson--Tropp (HMT) randomised SVD is benchmarked on
$N\times N$ float64 matrices at ranks $R\in\{8,16,32,64\}$ with
oversampling $p=10$.  The algorithm:
(1)~draw Gaussian sketch $\Omega \in \mathbb{R}^{N\times(R+p)}$;
(2)~form $Y = A\Omega$ and QR-orthogonalise;
(3)~project $B = Q^\top A$;
(4)~thin SVD of the small matrix $B$.
All BLAS/LAPACK calls dispatched to Apple Accelerate.

\subsection{Results: $128\times128$ matrices}

\begin{center}
\begin{tabular}{rrrrl}
\toprule
Rank $R$ & Mean (ms) & Std (ms) & Min (ms) & Target $\leq 10\,\text{ms}$ \\
\midrule
8  & 0.22 & 0.00 & 0.21 & \checkmark \\
16 & 0.35 & 0.05 & 0.33 & \checkmark \\
32 & 1.10 & 0.58 & 0.62 & \checkmark \\
64 & 2.37 & 1.17 & 1.39 & \checkmark \\
\midrule
Full SVD (reference) & 2.91 & --- & --- & --- \\
\bottomrule
\end{tabular}
\end{center}

\subsection{Results: $256\times256$ matrices}

\begin{center}
\begin{tabular}{rrrrl}
\toprule
Rank $R$ & Mean (ms) & Std (ms) & Min (ms) & Target $\leq 10\,\text{ms}$ \\
\midrule
8  & 0.80 & 0.38 & 0.43 & \checkmark \\
16 & 1.15 & 0.67 & 0.61 & \checkmark \\
32 & 2.09 & 1.27 & 1.13 & \checkmark \\
64 & 4.15 & 1.34 & 2.69 & \checkmark \\
\bottomrule
\end{tabular}
\end{center}

\subsection{Speedup over full SVD ($128\times128$)}

\begin{center}
\begin{tabular}{rrr}
\toprule
Rank $R$ & Rand SVD (ms) & Speedup \\
\midrule
8  & 0.22 & $13.5\times$ \\
16 & 0.35 & $8.4\times$ \\
32 & 1.10 & $2.7\times$ \\
64 & 2.37 & $1.2\times$ \\
\bottomrule
\end{tabular}
\end{center}

\begin{passbox}
\textbf{D4b: PASS.}
All $(N, R)$ pairs are within the $\leq 10\,\text{ms}$ target.
The critical $R=32$, $128\times128$ case gives $1.10\,\text{ms}$ mean,
a $2.7\times$ speedup over full SVD.
The high std at $R=32$ ($0.58\,\text{ms}$) reflects OS scheduling jitter
on the M1 CPU; the min of $0.62\,\text{ms}$ is the hardware floor.
\end{passbox}

% =============================================================================
\section{Memory Budget (Step 3/4 and 4/4)}
\label{sec:mem}
% =============================================================================

\subsection{Formulas}

\[
  M_\text{state}^\text{full}
  = N_x N_y N_z \cdot N_f \cdot 8\,\text{B},
  \qquad
  M_\text{peak}^\text{full}
  = 5 \times 1.2 \times M_\text{state}^\text{full}.
\]
\[
  M_\text{state}^\text{DLRA}
  = \bigl[(N_x + N_y + N_z)R + R^3\bigr] \cdot N_f \cdot 8\,\text{B},
  \qquad
  M_\text{peak}^\text{DLRA}
  = 6 \times M_\text{state}^\text{DLRA}.
\]

\subsection{Results}

\begin{center}
\begin{tabular}{lrrr}
\toprule
Configuration & State (GB) & Peak (GB) & $\leq 13\,\text{GB}$ \\
\midrule
$64\times64\times32$, full-rank   & 0.007 & 0.041 & \checkmark \\
$128\times128\times64$, full-rank & 0.055 & 0.328 & \checkmark \\
$128\times128\times64$, $R=32$    & 0.002 & 0.013 & \checkmark \\
$256\times256\times128$, $R=32$   & 0.003 & 0.017 & \checkmark \\
$512\times512\times128$, $R=32$   & 0.004 & 0.022 & \checkmark \\
\midrule
Hardware ceiling (operating) & \multicolumn{3}{l}{$13.0\,\text{GB}$} \\
Physical RAM & \multicolumn{3}{l}{$16.0\,\text{GB}$} \\
\bottomrule
\end{tabular}
\end{center}

\begin{passbox}
\textbf{D3: PASS.}
All configurations are within the 13\,GB operating ceiling.
Memory is not the bottleneck at any target resolution.
See \texttt{docs/phase0\_architecture.tex} \S3.3 for the note on the
proposal Table~5 discrepancy.
\end{passbox}

% =============================================================================
\section{Consolidated Summary and Go/No-Go}
\label{sec:summary}
% =============================================================================

\begin{center}
\begin{tabular}{llllll}
\toprule
ID & Check & Measured & Threshold & Status \\
\midrule
D1 & Parameter set sealed        & ---            & ---           &
  \textcolor{green!60!black}{\checkmark} \\
D2 & Benchmark thresholds sealed & ---            & ---           &
  \textcolor{green!60!black}{\checkmark} \\
D3 & Memory, R=32, 128³          & 0.013\,GB peak & $\leq 13$\,GB &
  \textcolor{green!60!black}{\checkmark} \\
D4a & Hero Run (128³, CPU)       & 0.41\,days     & $\leq 5$\,days &
  \textcolor{green!60!black}{\checkmark} \\
D4b & SVD, R=32, 128×128         & 1.10\,ms       & $\leq 10$\,ms  &
  \textcolor{green!60!black}{\checkmark} \\
\midrule
\textbf{Overall} & & & &
  \textcolor{green!60!black}{\textbf{ALL PASS}} \\
\bottomrule
\end{tabular}
\end{center}

\bigskip
\begin{passbox}
\textbf{Decision: Proceed to Phase~1 (3D spectral engine).}\\[4pt]
All four architecture decisions are resolved.
The parameter set, benchmark thresholds, memory budget, and hardware
performance are verified.
Phase~1 development begins with \texttt{src/engine/spectral.py}.

\medskip
\textbf{Before writing Phase~1 code, complete the three action items:}
\begin{enumerate}
  \item Commit \texttt{docs/phase0\_architecture.tex} (this file's companion)
    to the repository HEAD.
  \item Enable \texttt{jax\_enable\_x64 = True} in all simulation entry
    points.
  \item Switch development to the M4 16\,GB machine.
\end{enumerate}
\end{passbox}

% =============================================================================
\end{document}
% =============================================================================